\begin{statementbox}
	\textbf{\textcolor{CStatement}{条件}}
	\begin{description}[style=nextline, leftmargin=2.8em, labelsep=0.5em, itemsep=0.35em]
		\item[\textbf{\textcolor{CStatement}{条件 1（三点共线与等角）：}}]
			点 $A,B,C$ 共线，且 $\angle ABD=\angle BCE=\angle DBE=\theta$。
		\item[\textbf{\textcolor{CStatement}{条件 2（平行）：}}]
			直线 $AD\parallel BE$。
	\end{description}
	\textbf{\textcolor{CStatement}{结论}}
	\begin{description}[style=nextline, leftmargin=2.8em, labelsep=0.5em, itemsep=0.45em]
		\item[\textbf{\textcolor{CStatement}{结论一（相似与成比例）：}}]
			\textbf{\textcolor{CStatement}{直观描述：}}
			两个三角形相似，对应边成比例。
			\[
				\triangle ABD\sim\triangle BCE,\quad \frac{AB}{BC}=\frac{BD}{CE}=\frac{AD}{BE}.
			\]
		\item[\textbf{\textcolor{CStatement}{推广结论（全等）：}}]
			\textbf{\textcolor{CStatement}{直观描述：}}
			若有一组对应边相等，则两三角形全等。 当 $AB=BC$ 或 $BD=CE$ 时，有
			\[
				\triangle ABD\cong\triangle BCE.
			\]
		\item[\textbf{\textcolor{CStatement}{特例（一线三垂直）：}}]
			当 $\theta=90^{\circ}$ 时，条件 $AD\parallel BE$ 常可由垂直与共线推出，模型称为一线三垂直，常用于坐标系中构造全等或相似求点坐标。
	\end{description}

	\medskip
	\textbf{\textcolor{CStatement}{适用范围：}}
	平面几何中证明两三角形相似或全等；在平面直角坐标系中构造直角三角形（一线三垂直）求点坐标、最值及存在性问题。
\end{statementbox}
